% =============================================================================
%                              BRACED OPERATORS
% =============================================================================
%
% Braced operator support loaded by the braces package option.
%
% This module configures paired-delimiter star behavior, defines a template for
% operators that optionally accept bracketed, parenthesized, or braced operands,
% and upgrades standard AMS operators to the same interface.

% =============================================================================
%                            OPERATOR DELIMITERS
% =============================================================================

% -----------------------------------------------------------------------------
% User Commands
% -----------------------------------------------------------------------------
%
% \leftoperatorbrace and \rightoperatorbrace
%
% Purpose:
%   Provide mathcommand-aware left and right parentheses for operator operands.
%
% Usage:
%   Internal braced-operator rendering.
%
% \leftoperatorbracket and \rightoperatorbracket
%
% Purpose:
%   Provide mathcommand-aware left and right brackets for operator operands.
%
% Usage:
%   Internal braced-operator rendering.

\newmathcommand{\leftoperatorbrace}{\mlparen(}
\newmathcommand{\rightoperatorbrace}{\mrparen)}
\newmathcommand{\leftoperatorbracket}{\mlbrack[}
\newmathcommand{\rightoperatorbracket}{\mrbrack]}

% =============================================================================
%                        PAIRED DELIMITER CONFIGURATION
% =============================================================================

% -----------------------------------------------------------------------------
% Internal Functions
% -----------------------------------------------------------------------------
%
% \swapifbranches
%
% Purpose:
%   Reverses the true and false branches passed to a conditional. It is used to
%   invert mathtools' paired-delimiter star behavior so unstarred delimiters can
%   use the autoscaling branch.
%
% Parameters:
%   #1 - Conditional command to invoke.
%   #2 - Original true branch.
%   #3 - Original false branch.
%
% Usage:
%   Internal patch helper for \DeclarePairedDelimiter variants.

\newcommand\swapifbranches[3]{#1{#3}{#2}}

% -----------------------------------------------------------------------------
% Mathtools Patches
% -----------------------------------------------------------------------------
%
% Patch paired-delimiter declarations while mathtools' internal syntax is
% enabled. The braces option intentionally swaps starred and unstarred sizing
% behavior package-wide.

\MHInternalSyntaxOn
\patchcmd{\DeclarePairedDelimiter}{\@ifstar}{\swapifbranches\@ifstar}{}{}
\patchcmd{\DeclarePairedDelimiterXPP}{\@ifstar}{\swapifbranches\@ifstar}{}{}
\MHInternalSyntaxOff

% =============================================================================
%                         BRACED OPERATOR TEMPLATE
% =============================================================================

\ExplSyntaxOn

% -----------------------------------------------------------------------------
% Internal Templates
% -----------------------------------------------------------------------------
%
% \tmpl_braced_math_operator
%
% Purpose:
%   Renders an operator name, any absorbed PIE tokens, and an optional operand
%   enclosed in square brackets, parentheses, or braces.
%
% Template parameters:
%   <s> - Starred template instance uses \operatorname* for limit-style output.
%   <m> - Operator name text.
%
% Runtime parameters:
%   #1 - Star flag from the instantiated template.
%   #2 - Operator name supplied by the template instance.
%   #3 - Prime tokens absorbed by mathcommand.
%   #4 - Index tokens absorbed by mathcommand.
%   #5 - Exponent tokens absorbed by mathcommand.
%   #6 - Star flag disabling \mleft and \mright around operands.
%   #7 - Optional square-bracket operand.
%   #8 - Optional parenthesized operand.
%   #9 - Optional braced operand.
%
% Usage:
%   \TemplateInstancePIE{\rank}{\tmpl_braced_math_operator}<rank>

\NewTemplateCommand\tmpl_braced_math_operator<sm>{mmmsd[]d()g}{
  \IfBooleanTF{#1}{
    \cs_set:cpn {internal_operator_name:n} ##1 {\operatorname*{##1}}
  }{
    \cs_set:cpn {internal_operator_name:n} ##1 {\operatorname{##1}}
  }
  \IfBooleanTF{#6}{
    \cs_set:cpn {internal_l_brace_var:n} ##1 {##1}
    \cs_set:cpn {internal_r_brace_var:n} ##1 {##1}
  }{
    \cs_set:cpn {internal_l_brace_var:n} ##1 {\mleft##1}
    \cs_set:cpn {internal_r_brace_var:n} ##1 {\mright##1}
  }
  \IfNoValueTF{#7}{
    \IfNoValueTF{#8}{
      \IfNoValueTF{#9}{
        \internal_operator_name:n {#2} #3 #4 #5
      }{
        \internal_operator_name:n {#2} #3 #4 #5 \internal_l_brace_var:n\{ #9 \internal_r_brace_var:n\}
      }
    }{
      \internal_operator_name:n {#2} #3 #4 #5 \internal_l_brace_var:n( #8 \internal_r_brace_var:n) \IfNoValueTF{#9}{}{\bgroup #9 \egroup}
    }
  }{
    \internal_operator_name:n {#2} #3 #4 #5 \internal_l_brace_var:n[ #7 \internal_r_brace_var:n] \IfNoValueTF{#8}{}{(#8)} \IfNoValueTF{#9}{}{\bgroup #9 \egroup}
  }
}

% =============================================================================
%                          PUBLIC OPERATOR DECLARATION
% =============================================================================

% -----------------------------------------------------------------------------
% User Commands
% -----------------------------------------------------------------------------
%
% \DeclareBracedMathOperator
%
% Purpose:
%   Declares a PIE-aware operator that can optionally take bracketed,
%   parenthesized, or braced operands.
%
% Parameters:
%   *  - Optional star declares a limit-style operator with \operatorname*.
%   #1 - Operator command to define.
%   #2 - Operator name to typeset.
%
% Generated command usage:
%   \op
%   \op_i^2
%   \op[argument]
%   \op(argument)
%   \op{argument}

\RenewDocumentCommand \DeclareBracedMathOperator {smm} {
  \IfBooleanTF{#1}{
    \TemplateInstancePIE{#2}{\tmpl_braced_math_operator*}<#3>
  }{
    \TemplateInstancePIE{#2}{\tmpl_braced_math_operator}<#3>
  }
}

% =============================================================================
%                         STANDARD OPERATOR REDEFINITION
% =============================================================================

% -----------------------------------------------------------------------------
% Internal Functions
% -----------------------------------------------------------------------------
%
% \__xamsmath_redefine_operators:n
%
% Purpose:
%   Undefines an existing non-limit operator and redeclares it through
%   \DeclareBracedMathOperator.
%
% Parameters:
%   #1 - Operator name without a leading backslash.
%
% Usage:
%   Internal mapping over comma lists of AMS operator names.

\cs_new:Nn \__xamsmath_redefine_operators:n{
  \cs_undefine:c {#1}
  \exp_args:Nc \DeclareBracedMathOperator {#1} {#1}
}

% -----------------------------------------------------------------------------
% Internal Functions
% -----------------------------------------------------------------------------
%
% \__xamsmath_redefine_operators_star:n
%
% Purpose:
%   Undefines an existing limit-style operator and redeclares it through starred
%   \DeclareBracedMathOperator.
%
% Parameters:
%   #1 - Operator name without a leading backslash.
%
% Usage:
%   Internal mapping over comma lists of AMS operator names.

\cs_new:Nn \__xamsmath_redefine_operators_star:n{
  \cs_undefine:c {#1}
  \exp_args:NNc \DeclareBracedMathOperator* {#1} {#1}
}

% -----------------------------------------------------------------------------
% Operator Lists
% -----------------------------------------------------------------------------
%
% Standard AMS operator groups are stored by semantic category so each group can
% be upgraded through the correct starred or unstarred declaration helper.

\clist_set:Nn\__ams_trig_functions_list{arcsin,arctan,cos,cosh,cot,coth,csc,sec,sin,sinh,tan,tanh}
\clist_set:Nn\__ams_explog_functions_list{exp,log,ln,lg}
\clist_set:Nn\__ams_algebraic_functions_list{arg,deg,det,dim,gcd,lcm,ker,im,hom}
\clist_set:Nn\__ams_limit_operators_list{inf,sup,max,min,lin}

% -----------------------------------------------------------------------------
% Bulk Redefinition
% -----------------------------------------------------------------------------
%
% Non-limit groups are redeclared without a star. Limit operators are redeclared
% with a star so subscripts and superscripts behave like AMS limit operators.

\clist_map_function:NN\__ams_trig_functions_list\__xamsmath_redefine_operators:n
\clist_map_function:NN\__ams_explog_functions_list\__xamsmath_redefine_operators:n
\clist_map_function:NN\__ams_algebraic_functions_list\__xamsmath_redefine_operators:n
\clist_map_function:NN\__ams_limit_operators_list\__xamsmath_redefine_operators_star:n

% =============================================================================
%                           SPECIAL LIMIT OPERATORS
% =============================================================================

% -----------------------------------------------------------------------------
% User Commands
% -----------------------------------------------------------------------------
%
% The following AMS limit operators have printed names that do not match their
% control sequence names exactly. Each is undefined before being redeclared as a
% braced, limit-style operator.
%
% Usage:
%   \liminf_{n \to \infty} a_n
%   \injlim(A_i)
%   \varprojlim{X_i}

\cs_undefine:N \injlim
\DeclareBracedMathOperator* \injlim {inj \mkern 2mu lim}
\cs_undefine:N \projlim
\DeclareBracedMathOperator* \projlim {proj \mkern 2mu lim}
\cs_undefine:N \liminf
\DeclareBracedMathOperator* \liminf {lim \mkern 2mu inf}
\cs_undefine:N \limsup
\DeclareBracedMathOperator* \limsup {lim \mkern 2mu sup}
\cs_undefine:N \varliminf
\DeclareBracedMathOperator* \varliminf {\underline{lim}}
\cs_undefine:N \varlimsup
\DeclareBracedMathOperator* \varlimsup {\overline{lim}}
\cs_undefine:N \varinjlim
\DeclareBracedMathOperator* \varinjlim {\underrightarrow{\mathrm{lim}}}
\cs_undefine:N \varprojlim
\DeclareBracedMathOperator* \varprojlim {\underleftarrow{\mathrm{lim}}}

\ExplSyntaxOff
